\documentclass[11pt]{article}
\usepackage[a3paper, landscape, margin=1.8cm]{geometry}
\usepackage{multicol}
\usepackage{xcolor}
\usepackage{amsmath,amssymb}
\usepackage{tcolorbox}
\usepackage{graphicx}

\definecolor{foxorange}{RGB}{255, 107, 53}
\definecolor{nightblue}{RGB}{14, 20, 32}
\definecolor{mintgreen}{RGB}{16, 185, 129}
\definecolor{cardbg}{RGB}{250, 252, 255}

\tcbset{
  colback=cardbg,
  colframe=foxorange,
  fonttitle=\large\bfseries,
  coltitle=white,
  boxrule=2pt,
  arc=3mm,
  left=8mm, right=8mm, top=6mm, bottom=6mm
}

\begin{document}
\pagestyle{empty}

% ENCABEZADO
\begin{tcolorbox}[colback=nightblue, colframe=foxorange, boxrule=2.5pt, arc=4mm]
  \centering
  \vspace{0.2cm}
  {\Huge \bfseries \color{white} ANÁLISIS DE SISTEMAS ENERGÉTICOS Y EFICIENCIA TÉRMICA}\\[0.4cm]
  {\Large \color{foxorange} \textbf{Autor: Nombre del Investigador}} \quad 
  {\large \color{gray!30} | \quad Facultad de Ingeniería y Ciencias Aplicadas}\\[0.2cm]
  {\normalsize \color{gray!40} Asignatura: Termodinámica Avanzada \quad $\bullet$ \quad Profesor Guía: Dr. R. Silva \quad $\bullet$ \quad 2026}
  \vspace{0.2cm}
\end{tcolorbox}

\vspace{0.5cm}

\setlength{\columnsep}{1.6cm}
\begin{multicols}{2}

  \begin{tcolorbox}[title=1. Resumen y Justificación del Proyecto]
    La optimización de ciclos termodinámicos mediante fluidos orgánicos de bajo impacto ambiental representa una solución viable para la generación descentralizada de energía limpia.
    \begin{itemize}
      \item \textbf{Problema:} Pérdida de calor residual en plantas térmicas industriales.
      \item \textbf{Propuesta:} Ciclo Rankine Orgánico (ORC) con mezcla zeotrópica.
      \item \textbf{Meta:} Incrementar la eficiencia neta en más de un 12\%.
    \end{itemize}
  \end{tcolorbox}

  \vspace{0.4cm}

  \begin{tcolorbox}[title=2. Formulación y Balance de Energía]
    El balance de primera y segunda ley de la termodinámica para el evaporador y la turbina se define como:
    \begin{equation*}
      \eta_{\text{térmica}} = \frac{\dot{W}_{\text{turbina}} - \dot{W}_{\text{bomba}}}{\dot{Q}_{\text{in}}} = 1 - \frac{h_4 - h_1}{h_3 - h_2}
    \end{equation*}
    Donde $h_i$ son las entalpías específicas calculadas en cada estado de saturación.
  \end{tcolorbox}

  \vspace{0.4cm}

  \begin{tcolorbox}[title=3. Simulación y Curvas de Rendimiento]
    Se realizaron simulaciones iterativas variando la presión de condensación $P_c \in [1.2, 3.5]\text{ bar}$:
    \begin{center}
      \begin{tabular}{|l|c|c|c|}
        \hline
        \textbf{Configuración} & \textbf{Potencia (kW)} & \textbf{Eficiencia $\eta$} & \textbf{Emisiones} \\
        \hline
        Ciclo Tradicional & 45.2 & 14.3\% & Alta \\
        Flujo R134a & 58.7 & 18.9\% & Media \\
        \textbf{Propuesta Zeotrópica} & \textbf{72.4} & \textbf{23.1\%} & \textbf{Cero ODP} \\
        \hline
      \end{tabular}
    \end{center}
  \end{tcolorbox}

  \vspace{0.4cm}

  \begin{tcolorbox}[title=4. Conclusiones y Aplicabilidad]
    \begin{enumerate}
      \item La mezcla propuesta permite aprovechar fuentes térmicas de baja entalpía ($T < 140^\circ\text{C}$).
      \item El retorno de inversión estimado para una planta de mediana escala es inferior a 2.4 años.
    \end{enumerate}
    \vspace{0.2cm}
    \color{gray} \footnotesize Repositorio del modelo y código de simulación disponible en GitHub con licencia abierta.
  \end{tcolorbox}

\end{multicols}

\end{document}
